\documentclass[12pt,a4paper]{ctexart}

% 页面设置
\usepackage[a4paper,margin=2.5cm]{geometry}

% 数学环境
\usepackage{amsmath}
\usepackage{amssymb}

% 超链接与目录
\usepackage[hidelinks]{hyperref}

% 列表环境
\usepackage{enumitem}

% 段落设置
\setlength{\parindent}{2em}
\setlength{\parskip}{0.5em}

% 目录深度
\setcounter{tocdepth}{2}

\begin{document}

%------------------------------------------------
% 标题
%------------------------------------------------
\begin{titlepage}
    \centering
    \vspace*{4cm}

    {\Huge\bfseries 角色小传——王成\par}
    \vspace{2cm}

    {\Large XX高中校长\par}
    \vfill

    {\large 角色设定文档\par}
    \vspace{1cm}

    {\large \today\par}
\end{titlepage}

%------------------------------------------------
% 目录
%------------------------------------------------
\tableofcontents
\newpage

%================================================
% 第一部分
%================================================
\section{角色基本信息}

\begin{itemize}[leftmargin=2em]
    \item \textbf{姓名：}王成
    \item \textbf{年龄：}52岁
    \item \textbf{性别：}男
    \item \textbf{职业：}XX高中校长
\end{itemize}

王成是一名在学校工作了二十多年的教育管理者，最终成为XX高中的校长。他十分关心学生，经常把“学生综合素养”“全面发展”“先进教育理念”等词挂在嘴边，但他更加在意学校的排名、上级领导的评价以及自己的工作成绩。

他执行上级要求，并通过各种活动、数据、奖项和宣传材料证明学校取得了成绩。

%================================================
% 第二部分
%================================================
\section{人物弧光}

王成不思考什么是真正的教育，而只负责让学校看起来符合“先进教育理念”。

\[
    \text{普通管理者} \rightarrow \text{追求政绩} \rightarrow \text{形式主义}
    \rightarrow \text{机械执行} \rightarrow \text{教育机器}
\]

%================================================
% 第三部分
%================================================
\section{角色性格}

\subsection{表面和蔼，实际上强势}

王成讲话时始终面带微笑，经常使用“好不好？”“大家赞不赞同？”等表达，看起来是在征求意见，但实际上决定早已做出。

\subsection{重视成绩和学校形象}

王成并不真正关心兴趣班是否能够帮助学生成长，而更加关心兴趣班能否形成漂亮的宣传成果。

\subsection{善于察言观色}

面对不同的人，他会使用不同的话术。

面对老师，他强调“为了学生”；面对领导，他强调“教育成果”；面对媒体，他强调“学校取得了巨大成绩”。

\subsection{自洽}

王成最大的讽刺之处在于，他始终认为自己是在“提高学生综合素质”。

他并没有意识到，自己已经把学生变成宣传学校成果的工具。

%================================================
% 第四部分
%================================================
\section{角色口头禅}

王成最常说的一句话是：

\begin{quote}
    “为了学生的全面发展，好不好？”
\end{quote}

这句话通常出现在他宣布一个已经决定好的事情之后。

例如：

\begin{quote}
    “为了全面提高学生的综合素养，我们将星期六由放假改为上课，开设各种课外兴趣培训班，好不好？”
\end{quote}

随着剧情的发展，这句话可以逐渐变得机械化。

到了故事后期，即使面对完全不同的事情，王成依然会下意识地说出这句话，从而形成一种荒诞的重复。

%================================================
% 第五部分
%================================================
\section{角色习惯动作}

\begin{enumerate}[leftmargin=2.5em]
    \item \textbf{带头鼓掌三下：}
          每当王成说完一句“正确”的话，他会率先鼓掌三下，周围的人随后整齐地跟着鼓掌。

    \item \textbf{整理领带：}
          面对上级领导时，他会下意识整理领带和衣领，表现出紧张和讨好的状态。

    \item \textbf{听到“领导”立即挺直身体：}
          平时他的姿态比较放松，但一听到“领导来了”，身体会立刻挺直，脸上出现标准化的笑容。

    \item \textbf{笑声突然停止：}
          当领导停止笑或者表情发生变化时，王成会立即停止笑容，形成明显的机械化效果。

    \item \textbf{抬手控制现场：}
          当需要所有人鼓掌、安静或者微笑时，他都会轻轻抬手，像按下一个机械按钮。
\end{enumerate}

%================================================
% 第六部分
%================================================
\section{角色关键设计元素}

\subsection{服装}

王成始终穿着深色正式西装，领带系得非常整齐。

服装没有明显的个人特色，体现他逐渐成为整个管理系统的一部分。

\subsection{道具}

王成经常使用的道具包括扩音器、会议文件、手机、学校宣传册以及校长办公室里的奖杯。

其中，\textbf{宣传册和奖杯}是重要的视觉符号，代表王成最在意的“荣誉”“成绩”和“宣传成果”。

\subsection{色彩变化}

王成第一次出现时，可以使用相对正常、温和的色调。

随着剧情发展，画面逐渐变得冷硬、单调，使人物的“机器化”过程更加明显。

\subsection{身体语言}

王成的身体动作随着故事发展越来越标准化：

\[
    \text{说话} \rightarrow \text{微笑} \rightarrow \text{抬手}
    \rightarrow \text{鼓掌} \rightarrow \text{停止} \rightarrow \text{转身}
    \rightarrow \text{机械化行走}
\]

动作越整齐、越重复，人物的黑色幽默效果就越明显。

\subsection{核心视觉符号：机械按钮}

“抬手”是王成最重要的视觉动作。

\begin{itemize}[leftmargin=2em]
    \item 他抬手——所有人鼓掌。
    \item 他放手——所有人停止。
    \item 他微笑——所有人微笑。
    \item 他严肃——所有人安静。
\end{itemize}

这个动作让王成既像是“机器的一部分”，又像是在操纵其他人的机器。

%================================================
% 第七部分
%================================================
\section{角色核心矛盾}

王成最大的矛盾是：

\begin{quote}
    \textbf{他口口声声说自己是在培养学生，却越来越不关心学生真正需要什么。}
\end{quote}

他逐渐相信：

\begin{quote}
    “只要结果看起来正确，过程就不重要。”
\end{quote}

因此，王成与学生之间最核心的矛盾可以概括为：

\[
    \boxed{
        \text{学生的真实需求} \quad\mathrm{VS}\quad \text{学校需要展示的“教育成果”}
    }
\]

而王成始终选择后者。

%================================================
% 第八部分
%================================================
\section{角色总结}

王成不是传统意义上的反派。

他最大的特点不是“坏”，而是逐渐失去了判断什么真正重要的能力。

他把“提高综合素质”变成口号，把兴趣班变成政绩，把学生变成宣传材料，把掌声变成服从的证明。

最终，他自己也成为了这个荒诞教育系统中的一台“机器”。

\vspace{1em}

\begin{center}
    \textbf{人物关键词}

    \vspace{0.5em}

    校长、形式主义、政绩、服从、虚伪、机械化、教育异化、黑色幽默
\end{center}

\end{document}